% ============================================================================
% LANGENTWURF LE-018: Cloud-Speicher
% Profil: S (Senioren 65+) | Tier: 2 (Standard)
% ============================================================================

\documentclass[11pt,a4paper]{article}
\usepackage[utf8]{inputenc}
\usepackage[T1]{fontenc}
\usepackage[ngerman]{babel}
\usepackage{geometry}
\usepackage{fancyhdr}
\usepackage{lastpage}
\usepackage{xcolor}
\usepackage{tcolorbox}
\usepackage{enumitem}
\usepackage{longtable}
\usepackage{hyperref}
\usepackage{amssymb}
\usepackage{eurosym}

\geometry{left=2.5cm, right=2.5cm, top=2.5cm, bottom=2.5cm}
\definecolor{fach}{HTML}{b35c00}
\definecolor{gray}{HTML}{666666}
\definecolor{successgreen}{HTML}{2a7a4b}

\tcbuselibrary{skins,breakable}
\newtcolorbox{scorebox}{ colback=white, colframe=fach, fonttitle=\bfseries, title=Didaktik-Score, breakable }
\newtcolorbox{infobox}[1][]{ colback=fach!5, colframe=fach, fonttitle=\bfseries, title=#1, breakable }

\pagestyle{fancy}
\fancyhf{}
\fancyhead[L]{\small\textcolor{gray}{Digitale Grundlagen -- 018 Cloud-Speicher}}
\fancyhead[R]{\small\textcolor{gray}{LE Tier 2 / Profil S}}
\fancyfoot[C]{\small\thepage{} / \pageref{LastPage}}
\renewcommand{\headrulewidth}{0.4pt}

\begin{document}

\begin{titlepage}
    \centering
    \vspace*{2cm}
    {\Large\textcolor{gray}{Digitale Grundlagen -- Senioren 65+}}\\[2cm]
    {\Huge\textcolor{fach}{\textbf{Langentwurf}}}\\[0.5cm]
    {\large Tier 2 | Profil S}\\[2cm]
    {\LARGE\textbf{018 -- Cloud-Speicher verstehen}}\\[0.5cm]
    {\large Modul D: Produktivität}\\[2cm]
    \begin{tabular}{rl}
        \textbf{Zeitrahmen:} & 60--75 Minuten \\
        \textbf{Vorkenntnisse:} & Fotos \& Dateien (017) \\
    \end{tabular}
    \vfill
    {\normalsize Stand: \today}
\end{titlepage}

\tableofcontents
\newpage

\section{Bedingungsanalyse}

\subsection{Ausgangssituation}
\begin{itemize}
    \item ``Cloud'' ist ein Mysterium
    \item ``Wo sind meine Daten eigentlich?''
    \item Angst vor Datenverlust und Fremdzugriff
    \item iCloud-Meldungen werden ignoriert oder verwirren
    \item 5 GB kostenlos sind schnell voll
\end{itemize}

\subsection{Typische Probleme}
\begin{itemize}
    \item ``Was ist die Cloud?''
    \item ``Sind meine Fotos sicher?''
    \item ``iCloud-Speicher voll -- was nun?''
    \item ``Muss ich dafür bezahlen?''
\end{itemize}

\section{Sachanalyse}

\subsection{Kernkonzepte}
\begin{enumerate}
    \item \textbf{Cloud:} Speicherplatz im Internet (auf Computern von Apple/Google)
    \item \textbf{iCloud:} Apples Cloud-Dienst (bei Apple-ID dabei)
    \item \textbf{Sync:} Automatische Kopie auf allen Geräten
    \item \textbf{Backup:} Sicherungskopie aller Daten
    \item \textbf{Speicherplan:} Kostenlos 5 GB, mehr gegen Gebühr
\end{enumerate}

\subsection{Alltagsanalogien}
\begin{tabular}{|l|p{8cm}|}
\hline
\textbf{Begriff} & \textbf{Analogie} \\
\hline
Cloud & Schließfach bei der Bank -- nicht bei dir zuhause, aber sicher \\
Sync & Automatische Kopie in dein Schließfach \\
Backup & Komplette Kopie aller Dokumente als Sicherheit \\
5 GB & Kleine Aktentasche -- passt nicht viel rein \\
50 GB & Großer Aktenkoffer -- mehr Platz \\
\hline
\end{tabular}

\section{Didaktische Analyse}

\subsection{Stellung in der Reihe}
Konzeptuelle Einheit. Wichtig für Verständnis, warum Daten ``überall'' sind.

\subsection{Didaktische Reduktion}
\textbf{Ausgeklammert:} Andere Cloud-Dienste (Google, Dropbox), technische Details, Familienfreigabe.

\textbf{Fokus:} iCloud verstehen, Speicher prüfen, einfaches Backup.

\begin{infobox}[Warum Cloud?]
\begin{itemize}
    \item \textbf{Sicherheit:} Handy verloren? Daten sind noch da!
    \item \textbf{Bequemlichkeit:} Foto auf iPhone gemacht, auf iPad anschauen
    \item \textbf{Speicherplatz:} Mehr Platz auf dem Gerät
\end{itemize}
\end{infobox}

\section{Lernziele}

\begin{tabular}{|c|p{7cm}|c|c|}
\hline
\textbf{Nr.} & \textbf{Die Teilnehmer können...} & \textbf{Stufe} & \textbf{Niveau} \\
\hline
FZ1 & erklären, was ``Cloud'' bedeutet & Verstehen & Basis \\
FZ2 & ihren iCloud-Speicherstand prüfen & Anwenden & Basis \\
FZ3 & verstehen, was in iCloud gespeichert wird & Verstehen & Basis \\
FZ4 & die Vorteile von Cloud-Speicher nennen & Verstehen & Basis \\
FZ5 & entscheiden, ob mehr Speicher sinnvoll ist & Anwenden & Vertiefung \\
\hline
\end{tabular}

\section{Verlaufsplanung}

\begin{longtable}{|p{1cm}|p{2cm}|p{5cm}|p{3cm}|p{2cm}|}
\hline
\textbf{Zeit} & \textbf{Phase} & \textbf{Inhalt} & \textbf{TN-Aktivität} & \textbf{Material} \\
\hline
0--15 & Einstieg & ``Was passiert, wenn Ihr Handy kaputt geht?'' Cloud-Konzept & Nachdenken & -- \\
\hline
15--30 & Was ist Cloud? & Schließfach-Analogie, iCloud erklären & Zuhören, Fragen & S.1 \\
\hline
30--45 & Speicher prüfen & Einstellungen $\rightarrow$ Apple-ID $\rightarrow$ iCloud & Eigenes Gerät & S.1 \\
\hline
45--60 & Was wird gespeichert? & Fotos, Notizen, Kontakte, Backup & Einstellungen ansehen & S.2 \\
\hline
60--70 & Speicher voll? & Optionen: Löschen, Upgrade, Was brauche ich? & Diskutieren & S.2-3 \\
\hline
\end{longtable}

\section{Lösungen}

\textbf{iCloud-Speicher prüfen:}
\begin{itemize}
    \item Einstellungen $\rightarrow$ [Ihr Name] oben $\rightarrow$ iCloud
    \item Farbiger Balken zeigt: Was verbraucht Speicher?
    \item Grün = Fotos, Gelb = Backups, Grau = Andere
\end{itemize}

\textbf{Was wird automatisch gespeichert?}
\begin{itemize}
    \item Fotos (wenn aktiviert)
    \item Kontakte, Kalender, Erinnerungen
    \item Notizen
    \item Safari-Lesezeichen
    \item Backup des gesamten Geräts
\end{itemize}

\textbf{Speicherplatz-Optionen:}
\begin{itemize}
    \item \textbf{5 GB kostenlos:} Reicht für Kontakte, Notizen, kleine Backups
    \item \textbf{50 GB (0,99\euro/Monat):} Genug für die meisten Nutzer
    \item \textbf{200 GB (2,99\euro/Monat):} Für viele Fotos und Videos
\end{itemize}

\textbf{Speicher sparen:}
\begin{itemize}
    \item Alte Backups anderer Geräte löschen
    \item Nicht alle Apps in iCloud sichern
    \item Fotos aufräumen (siehe 017)
\end{itemize}

\section{Didaktik-Score}

\begin{scorebox}
\begin{tabular}{|c|l|c|c|p{3.5cm}|}
\hline
\# & Dimension & Gew. & Score & Notiz \\
\hline
1 & Differenzierung & 15\% & 7/10 & Eher konzeptuelle Einheit \\
2 & Taxonomie & 10\% & 9/10 & Verstehen + Anwenden \\
3 & Alltagsrelevanz & 10\% & 9/10 & Datensicherheit wichtig \\
4 & Aktivierung & 10\% & 8/10 & Eigene Einstellungen prüfen \\
5 & Scaffolding & 10\% & 9/10 & Schließfach-Analogie \\
6 & Feedback & 10\% & 8/10 & Checkliste \\
7 & Sprachliche Klarheit & 10\% & 9/10 & Gute Erklärungen \\
8 & Motivation & 10\% & 9/10 & Sicherheit motiviert \\
9 & Barrierefreiheit & 10\% & 9/10 & Klare Pfade \\
10 & Praktikabilität & 5\% & 9/10 & 70 Min angemessen \\
\hline
\multicolumn{3}{|r|}{\textbf{GESAMT:}} & \textbf{8.65/10} & Knapp unter Freigabe \\
\hline
\end{tabular}
\end{scorebox}

\end{document}
